% Source: https://artofproblemsolving.com/community/c3h2503429_general_form_of_the_age_problem
% Stable identity: aops-topic-2503429-q01
\textbf{(a)} Let the present ages be $A$ and $B$. Then $A=xB$ and $A+y=z(B+y)$. Substituting,
\[xB+y=zB+zy\iff(x-z)B=(z-1)y,\]
so $B=\dfrac{y(z-1)}{x-z}$ and $A=xB=\dfrac{xy(z-1)}{x-z}$.

\textbf{(b)} With $y>0$, $B>0$ requires $z-1$ and $x-z$ to have the same sign. If both were negative, then $x<z<1$, contradicting $x>1$. So both are positive: $1<z<x$. (The ratio of the ages moves towards $1$ as time passes.)

\textbf{(c)} Here $x=4$, $z=2$, $y=20$, so the daughter is $\dfrac{20\cdot1}{2}=10$ and the mother is $40$. Check: in $20$ years they will be $30$ and $60$. So the ages are $\boxed{40\text{ and }10}$.
